\color{unimain}{\subsection{Фиксация положения выключателей}}
\color{black}

\begin{enumerate}[label=\arabic{section}.\arabic{subsection}.\arabic*,  align=left, labelsep=4pt, leftmargin=0pt, itemindent=57pt]

\item
В устройстве реализована функция фиксации положения выключателей В1 и В2 высокой стороны
трансформатора.  

\begin{figure}[H]
\centering
\includegraphics[width=0.9\textwidth,height=0.9\textheight,keepaspectratio]{\fbpath/761.3381 ФПВ/img/FPV.pdf}
\captionof{figure}{Функциональная схема ФПВ В1 и В2 ВН}\label{pict:FPV}
\end{figure}

\item
Функциональный блок фиксации положения выключателей (ФПВ) предназначен для определения
состояния выключателей на основании технологической информации:
\begin{itemize}
\item состояния блок-контактов выключателя;
\item положения разъединителей в цепи выключателя;
\item информации о выводе в ремонт выключателя.
\end{itemize}

\item
Контроль положения разъединителей при формировании сигналов включенного или отключенного
положения выключателей вводится программным переключателями \textbf{«Контр Р В1 ВН»} и \textbf{«Контр Р В2 ВН»}.
Также имеется возможность выбрать количество разъединителей, участвующих в формировании сигналов
включенного или отключенного положения выключателей, задается программными переключателями \textbf{«Схема
Р В1 ВН»} и \textbf{«Схема Р В2 ВН»}. А для выключателей В2 ВН имеется возможность выбрать тип выключателя
линейный или обходной, задается программным переключателем \textbf{«Тип В2 ВН»}. 

\item
Также функция ФПВ пофазно контролирует состояние блок-контактов выключателя и в случае
расхождения полюсов формирует сигналы \textit{«Внутр. пуск. ЗНФ В1 ВН»} и \textit{«Внутр. пуск. ЗНФ В2 ВН»}.


\small
\setlength{\extrarowheight}{0.05cm} % высота ячеек в таблице
\begin{longtable}
{|>{\centering\arraybackslash}m{4.5cm}|>{\centering\arraybackslash}m{4.5cm}|>{\centering\arraybackslash}m{3cm}|>{\centering\arraybackslash}m{3.5cm}|}
\caption{Перечень уставочных параметров ФПВ ВН\hfill\vspace{-0.5\baselineskip}}\label{tzt:tbl1} \\ 
\hhline{|-|-|-|-|}
\rowcolor{gray!30}
Наименование уставки & Диапазон & Шаг &  Значение по умолчанию \\ 
\hhline{|-|-|-|-|}
\endfirsthead

\caption*{\hspace{3pt}\emph{Продолжение таблицы \ref{tzt:tbl1}\hfill\vspace{-0.5\baselineskip}}} \\ % сделано по ГОСТ 2.105 п.6.8.7
\hhline{|-|-|-|-|}
\rowcolor{gray!30}
Наименование уставки & Диапазон & Шаг &  Значение по умолчанию \\  
\hhline{|-|-|-|-|}
\endhead
\hhline{|-|-|-|-|}

\endfoot
\endlastfoot

%%%%%%%%%%%%содержимое%%%%%%%%%%%%%%

\end{longtable}
\setlength{\extrarowheight}{0.0cm} % возврат параметра высота ячеек в таблице
\normalsize

\end{enumerate} 